% ============================================================
\chapter{Machine Learning in Finance}
\label{ch:ml_finance}
% ============================================================

\section{Introduction}

Machine learning\index{machine learning!finance} (ML) is transforming
quantitative finance by enabling the discovery of nonlinear structures in
market data. This chapter presents the main applications: algorithmic trading,
neural network option pricing, reinforcement learning for portfolio management,
NLP-based sentiment analysis, and associated risks and regulations.

\begin{intuition}
Classical models (Black-Scholes, GARCH) impose parametric structure. ML allows
the data to ``speak for itself'' by learning complex relationships directly,
at the cost of increased overfitting risk and reduced interpretability.
\end{intuition}

\section{Algorithmic Trading}

\begin{definition}[Algorithmic trading strategy]
\index{algorithmic trading}
An \textbf{algorithmic trading strategy} is a systematic rule
$\delta_t = g(x_t; \theta)$ determining the position $\delta_t$ at time $t$
as a function of observable features $x_t$ and parameters $\theta$ optimised
on historical data.
\end{definition}

\begin{proposition}[Classical features]
Typical predictors include:
\begin{enumerate}
    \item \textbf{Technical indicators}: moving averages, RSI, Bollinger bands,
    MACD.
    \item \textbf{Microstructure}: order flow imbalance, bid-ask spread, volume.
    \item \textbf{Fundamental features}: financial ratios, earnings surprises.
    \item \textbf{Alternative data}: satellite imagery, web traffic, social
    media sentiment.
\end{enumerate}
\end{proposition}

\begin{attention}[title=\textbf{Overfitting and data snooping}]
Extensive backtesting on historical data produces strategies that perform
perfectly in-sample but fail out-of-sample. Remedies include: time-series
cross-validation (walk-forward), Bonferroni correction for multiple testing,
and holding out fresh data for the final test.
\end{attention}

\begin{theorem}[Generalisation bound]
\index{generalisation bound}
For a hypothesis class $\mathcal{H}$ of VC dimension $d_{\mathrm{VC}}$, with
probability at least $1 - \delta$ on a sample of size $n$:
\[
    R(h) \leq \hat{R}(h) + \sqrt{\frac{d_{\mathrm{VC}}(\ln(2n/d_{\mathrm{VC}}) + 1)
    + \ln(4/\delta)}{n}},
\]
where $R(h)$ is the true risk and $\hat{R}(h)$ is the empirical risk.
\end{theorem}

\section{Neural Networks for Option Pricing}

\begin{definition}[Feedforward neural network]
\index{neural network}
A \textbf{feedforward neural network} with $L$ layers is a function
\[
    f_\theta(x) = W_L\,\sigma(W_{L-1}\,\sigma(\cdots\sigma(W_1 x + b_1)\cdots)
    + b_{L-1}) + b_L,
\]
where $\sigma$ is a nonlinear activation function (ReLU, tanh, etc.),
$W_\ell$ are weight matrices, and $b_\ell$ are biases.
\end{definition}

\begin{theorem}[Universal approximation]
\index{universal approximation}
For any continuous function $g : \R^d \to \R$ on a compact set $K$ and any
$\epsilon > 0$, there exists a single-hidden-layer feedforward network with
sigmoid activation such that
$\sup_{x \in K}\abs{f_\theta(x) - g(x)} < \epsilon$.
\end{theorem}

\begin{proposition}[Neural network pricing]
The approach consists of:
\begin{enumerate}
    \item Generate $N$ parameter sets $(S_0^{(i)}, K^{(i)}, T^{(i)},
    \sigma^{(i)}, r^{(i)})$ and compute prices $V^{(i)}$ via Black-Scholes
    or Monte Carlo.
    \item Train a network $f_\theta$ minimising
    $\frac{1}{N}\sum_{i=1}^N (f_\theta(x^{(i)}) - V^{(i)})^2$.
    \item At inference, $f_\theta$ yields prices and Greeks
    ($\Delta = \partial f_\theta/\partial S$) near-instantaneously.
\end{enumerate}
\end{proposition}

\begin{remark}
Neural networks accelerate pricing by a factor of $10^3$ to $10^5$ compared
to Monte Carlo, which is crucial for XVA computation (CVA, DVA, FVA) on
portfolios containing thousands of trades.
\end{remark}

\section{Reinforcement Learning for Portfolio Management}

\begin{definition}[Markov Decision Process (MDP)]
\index{MDP}
An \textbf{MDP} is a tuple $(\mathcal{S}, \mathcal{A}, P, R, \gamma)$ where
$\mathcal{S}$ is the state space, $\mathcal{A}$ is the action space,
$P(s'|s, a)$ is the transition dynamics, $R(s, a)$ is the reward, and
$\gamma \in [0, 1)$ is the discount factor.
\end{definition}

\begin{definition}[Policy and value function]
\index{policy!reinforcement learning}
A \textbf{policy} $\pi : \mathcal{S} \to \mathcal{A}$ specifies the action at
each state. The \textbf{value function} is
$V^\pi(s) = \E_\pi\!\left[\sum_{t=0}^\infty \gamma^t R(s_t, a_t) \mid
s_0 = s\right]$. The optimal policy satisfies
$V^*(s) = \max_\pi V^\pi(s)$.
\end{definition}

\begin{theorem}[Bellman equation]
\index{Bellman equation}
The optimal value function satisfies
\[
    V^*(s) = \max_{a \in \mathcal{A}} \left[R(s, a)
    + \gamma \sum_{s'} P(s'|s, a)\,V^*(s')\right].
\]
\end{theorem}

\begin{proposition}[Application to portfolio management]
\begin{itemize}
    \item \textbf{State}: $s_t = (w_t, S_t, \text{features}_t)$ (wealth,
    prices, indicators).
    \item \textbf{Action}: $a_t = (\delta_1, \ldots, \delta_d)$ (allocations
    across $d$ assets).
    \item \textbf{Reward}: $R_t = \ln(w_{t+1}/w_t)$ (log return) or
    differential Sharpe ratio.
    \item \textbf{Algorithms}: DDPG, PPO, SAC for continuous action spaces.
\end{itemize}
\end{proposition}

\section{NLP and Sentiment Analysis}

\begin{definition}[Financial sentiment analysis]
\index{sentiment analysis}
\textbf{Sentiment analysis} in finance consists of extracting a
bullish/bearish signal from textual sources: press releases, analyst reports,
tweets, earnings call transcripts.
\end{definition}

\begin{proposition}[NLP pipeline for finance]
\begin{enumerate}
    \item \textbf{Collection}: Reuters/Bloomberg feeds, Twitter/Reddit APIs.
    \item \textbf{Preprocessing}: tokenisation, stop word removal,
    lemmatisation.
    \item \textbf{Model}: FinBERT (BERT fine-tuned on financial text) or
    specialised GPT.
    \item \textbf{Aggregation}: weighted average sentiment score by source
    and recency.
    \item \textbf{Trading signal}: combination with quantitative features.
\end{enumerate}
\end{proposition}

\begin{keyformulas}[title=\textbf{ML Strategy Performance Metrics}]
\begin{align*}
    \text{Sharpe Ratio} &= \frac{\E[R_p] - r_f}{\sqrt{\Var(R_p)}} \\
    \text{Sortino Ratio} &= \frac{\E[R_p] - r_f}{\sqrt{\E[\min(R_p - r_f, 0)^2]}} \\
    \text{Maximum Drawdown} &= \max_{t}\frac{\max_{s \leq t} W_s - W_t}{\max_{s \leq t} W_s}
\end{align*}
\end{keyformulas}

\section{Risks and Regulation}

\begin{definition}[Model risk]
\index{model risk}
\textbf{Model risk} in ML includes:
\begin{itemize}
    \item \textbf{Overfitting}: the model memorises noise.
    \item \textbf{Distribution shift}: future data differs from training data
    (regime changes).
    \item \textbf{Black box}: inability to explain decisions
    (interpretability problem).
    \item \textbf{Adversarial attacks}: deliberate manipulation of inputs to
    deceive the model.
\end{itemize}
\end{definition}

\begin{remark}
The EU AI Act (2024) classifies credit scoring and algorithmic trading systems
as ``high risk,'' imposing transparency, documentation, and human oversight
requirements.
\end{remark}

\begin{codeblock}[title=\textbf{Option pricing with a neural network}]
\begin{minted}{python}
import numpy as np
import torch
import torch.nn as nn
from scipy.stats import norm

def bs_call(S, K, T, r, sigma):
    d1 = (np.log(S/K) + (r + sigma**2/2)*T) / (sigma*np.sqrt(T))
    d2 = d1 - sigma*np.sqrt(T)
    return S*norm.cdf(d1) - K*np.exp(-r*T)*norm.cdf(d2)

n_train = 50_000
S = np.random.uniform(50, 150, n_train)
K = np.random.uniform(50, 150, n_train)
T = np.random.uniform(0.1, 2.0, n_train)
sigma = np.random.uniform(0.1, 0.5, n_train)
prices = bs_call(S, K, T, 0.05, sigma)

X = torch.tensor(np.column_stack([S, K, T, sigma]), dtype=torch.float32)
y = torch.tensor(prices, dtype=torch.float32).unsqueeze(1)

model = nn.Sequential(
    nn.Linear(4, 64), nn.ReLU(),
    nn.Linear(64, 64), nn.ReLU(),
    nn.Linear(64, 1))
optimizer = torch.optim.Adam(model.parameters(), lr=1e-3)

for epoch in range(200):
    pred = model(X)
    loss = nn.MSELoss()(pred, y)
    optimizer.zero_grad(); loss.backward(); optimizer.step()
print(f"Final MSE: {loss.item():.6f}")
\end{minted}
\end{codeblock}

\section{Exercises}

\begin{exercise}[Feature engineering]
For an S\&P 500 trading strategy, propose 10 relevant features. Discuss
look-ahead bias and multicollinearity issues.
\end{exercise}

\begin{exercise}[Greeks via autodiff]
Train a neural network to price European calls. Compute delta via automatic
differentiation and compare with the analytical formula.
\end{exercise}

\begin{exercise}[RL for allocation]
Implement a simple DQN agent for allocation between a risky asset (return
$\mathcal{N}(8\%, 20\%)$) and a risk-free asset (3\%). Compare the learned
policy with Merton's rule.
\end{exercise}

\begin{exercise}[Sentiment analysis]
Collect 1000 financial news headlines and classify them as
positive/negative/neutral. Measure the correlation between the aggregated
sentiment score and next-day returns.
\end{exercise}

\begin{exercise}[Temporal validation]
Explain why standard k-fold cross-validation is inappropriate in finance.
Implement walk-forward validation with a 2-year training window and a 6-month
test window.
\end{exercise}

\begin{exercise}[Interpretability]
Apply Shapley values (SHAP) to an XGBoost model for return prediction.
Identify the most important features and discuss the stability of
explanations over time.
\end{exercise}
